\chapter{The APIx Specification}

This chapter is the formal definition of the index. It should be readable as a
specification document by someone who has never seen the code, and every step should be
implementable and testable in isolation. Treat it as the contract between the statistics
work and the engineering work.

\section{The product key: what exactly is being priced}
\label{sec:key}

A price index requires a \emph{constant-quality} product whose price is tracked over time.
For airfares this definition is the whole game --- get it wrong and every downstream number
is meaningless. APIx defines the priced product as:

\begin{keybox}
\centering
\bhead{One economy-class, one-way, one-adult, hand-baggage-inclusive seat on a
non-stop scheduled flight, on route $r$, operated by carrier $c$, departing on date $d$,
quoted at exactly $\tau$ days before departure, at a fixed collection clock time,
from an unauthenticated India-egress session.}
\end{keybox}

\noindent Each qualifier is load-bearing:

\begin{center}
\small
\begin{tabularx}{\textwidth}{@{}L{3.4cm} X@{}}
\toprule
\bhead{Qualifier} & \bhead{Why it is in the definition} \\
\midrule
Economy class & GST differs by cabin (5\% economy vs.\ 12\% other than economy --- \textbf{[VERIFY]}); mixing cabins mixes tax regimes and quality. \\
One-way & Return fares are bundled products with their own pricing logic; a separate return-basis index is a documented extension. \\
One adult & Child/infant fares follow different rules; multi-pax triggers group pricing. \\
Hand baggage only & The minimum common bundle across LCC and FSC. Checked bag is an ancillary and is indexed separately. \\
Non-stop & A one-stop itinerary is a different service (longer duration = lower quality). Mixing them is unadjusted quality change. \\
Fixed $\tau$ & The advance-purchase window \emph{is} a product characteristic. A T$+$7 fare and a T$+$45 fare are different products, not two observations of one. \\
Fixed clock time & Fares move intraday. Collecting at 02:00 IST one day and 14:00 the next injects diurnal noise into the price signal. \\
Unauthenticated, cookie-clean & Removes personalisation and loyalty pricing. A logged-out session is also the legally cleanest one (Part IV). \\
India egress & Currency, tax treatment and price display are geography-dependent. \\
\bottomrule
\end{tabularx}
\end{center}

\subsection{The matching key}

Two quotes are \emph{matched} (i.e.\ the same product in two periods) if and only if they
agree on the tuple
\begin{equation}
\kappa \;=\; \big(r,\; c,\; \tau,\; \text{flight\_number},\; \text{fare\_family } k,\; \text{cabin},\; \text{stops}\big)
\end{equation}
Matching on \code{flight\_number} rather than on departure date is deliberate: it tracks
\emph{the same flight in the schedule} (e.g.\ 6E-2134, the 07:10 \rt{del}$\to$\rt{bom})
across weeks, which is the closest airline analogue of tracking the same SKU in a shop. It
also survives the daily rotation of departure dates that $\tau$ being fixed necessarily
implies.

\begin{warnbox}
\bhead{Day-of-week is a confound and must be controlled.} If $\tau$ is fixed and $t$
advances one day, then $d$ advances one day too --- so a Tuesday-departure observation
becomes a Wednesday-departure observation. Since fares are strongly day-of-week dependent,
a naive daily index would report pure day-of-week variation as inflation.

APIx handles this in two ways: (i) the daily index is computed on a \textbf{7-day centred
geometric moving average} of the daily cell relatives, which annihilates a weekly cycle
exactly; and (ii) a day-of-week fixed effect is included in the TPD specification. State
this unprompted --- it is the kind of trap a MoSPI statistician will probe for, and having
the answer ready is worth more than any feature.
\end{warnbox}

\section{Stage 0: raw quote to cleaned price}

\subsection{Fare decomposition}

A displayed ``total fare'' on an Indian OTA is a sum of economically distinct components
with completely different inflation dynamics. Indexing the total without decomposing it
means a UDF revision at one airport shows up as ``airfare inflation''. The decomposition:

\begin{equation}
P^{\text{total}} \;=\; \underbrace{P^{\text{base}}}_{\text{airline, dynamic}}
\;+\; \underbrace{\text{YQ} + \text{YR}}_{\text{carrier surcharges}}
\;+\; \underbrace{\text{UDF} + \text{ASF} + \text{RCS}}_{\text{statutory / airport}}
\;+\; \underbrace{\text{GST}}_{\text{tax}}
\;+\; \underbrace{\text{CF}}_{\text{OTA convenience fee}}
\end{equation}

\begin{center}
\small
\begin{tabularx}{\textwidth}{@{}L{2.1cm} L{4.4cm} X@{}}
\toprule
\bhead{Component} & \bhead{Set by} & \bhead{Treatment in APIx} \\
\midrule
$P^{\text{base}}$ & Airline revenue management & \textbf{The core signal.} Primary index is built on this. \\
YQ / YR & Airline (fuel / carrier misc.) & Economically part of the fare; folded into $P^{\text{base}}$ for the headline index, reported separately in the diagnostic view. \\
UDF & \textbf{AERA} tariff orders, per airport & Administered price. Changes are step functions on known dates. Indexed separately. \\
ASF & \textbf{BCAS / MoCA} notification & Administered, uniform. Same treatment. \\
RCS levy & MoCA (UDAN cross-subsidy) & Applies on specified routes. Same treatment. \\
GST & GST Council; 5\% economy / 12\% other \textbf{[VERIFY]} & Ad valorem on the fare base. \\
CF & The OTA & \textbf{Not a fare.} It is a distribution charge that differs across OTAs for an identical seat. Excluded from the fare index; published as a separate ``Distribution Cost'' series. \\
\bottomrule
\end{tabularx}
\end{center}

\noindent APIx publishes \textbf{three headline variants}, which is itself a USP because it
lets a user see exactly how much of ``airfare inflation'' is airline behaviour versus
administered charges:

\begin{align}
\text{APIx-B} &: \text{base fare} + \text{YQ/YR only} && \text{(pure airline pricing)}\\
\text{APIx-T} &: \text{everything except OTA convenience fee} && \text{(what a traveller pays --- headline)}\\
\text{APIx-A} &: \text{APIx-T} + \text{OTA convenience fee} && \text{(all-in checkout cost)}
\end{align}

\begin{keybox}
\bhead{Do not scrape administered charges.} UDF, ASF, RCS and GST rates are published by
AERA, BCAS/MoCA and the GST Council. Maintain them as a \textbf{slowly-changing dimension
table} loaded from official notifications, versioned with effective-from/effective-to dates.
Then compute the base fare as a residual and cross-check it against the airline's own fare
breakup where one is displayed. This turns a scraping problem into a reconciliation problem
--- and reconciliation is auditable.
\end{keybox}

\subsection{Outlier detection}

Fare quotes contain genuine extremes (last-seat pricing) and genuine errors (parse
failures, currency confusion, promotional glitches). The cleaning rule must remove the
second without deleting the first, because the first is the signal.

Work in log space throughout: $y_i = \ln \tilde p_i$.

\begin{enumerate}
\item \textbf{Hard validity gates} (deterministic, not statistical). Reject if:
$\tilde p \le 0$; $\tilde p <$ the sum of known statutory components; currency $\neq$ INR;
$\tau$ mismatch $>0$ days; duration implausible for the great-circle distance; flight number
not in the schedule reference.

\item \textbf{Tukey fences within cell}, on logs. With $Q_1, Q_3$ the cell quartiles and
$\text{IQR}=Q_3-Q_1$:
\begin{equation}
\text{flag if } y_i < Q_1 - 3\,\text{IQR} \quad\text{or}\quad y_i > Q_3 + 3\,\text{IQR}
\end{equation}
The multiplier is $3$, not the usual $1.5$: fares are legitimately dispersed and $1.5$
would trim real observations.

\item \textbf{Hampel / MAD filter} on the \emph{price relative}, which is more robust than
quartiles for small cells. With $m$ the median and
$\text{MAD} = \operatorname{median}_i |y_i - m|$:
\begin{equation}
z_i^{\text{rob}} \;=\; \frac{0.6745\,(y_i - m)}{\text{MAD}},
\qquad \text{flag if } |z_i^{\text{rob}}| > 3.5
\end{equation}
(The constant $0.6745$ makes MAD a consistent estimator of $\sigma$ under normality.)

\item \textbf{Longitudinal jump test.} Flag a quote whose own log-relative to its previous
observation exceeds a cell-specific threshold: $|\ln(p_{i,t}/p_{i,t-1})| > 3\hat\sigma_g$.
This catches parse regressions that survive the cross-sectional filters.

\item \textbf{Disposition, not deletion.} A flagged quote is \textbf{never silently
dropped}. It is written to the store with \code{quality\_flag} set and an explicit
\code{disposition} in \{\code{ACCEPTED}, \code{WINSORISED}, \code{EXCLUDED},
\code{QUARANTINED}\}. The published index uses \code{ACCEPTED} $\cup$ \code{WINSORISED};
the audit endpoint exposes all four. Winsorisation caps at the fence rather than removing:
\begin{equation}
y_i^{w} = \min\big(\max(y_i,\; Q_1 - 3\,\text{IQR}),\; Q_3 + 3\,\text{IQR}\big)
\end{equation}
\end{enumerate}

\subsection{Imputation for missing quotes}

If a cell has a matched product that yields no quote at $t$ (page failed, flight cancelled,
inventory sold out), leaving a gap biases the index. Three cases, three treatments:

\begin{center}
\small
\begin{tabularx}{\textwidth}{@{}L{3.6cm} L{4.2cm} X@{}}
\toprule
\bhead{Cause} & \bhead{Detectable by} & \bhead{Treatment} \\
\midrule
Collection failure & HTTP error, empty parse, circuit-breaker trip & \textbf{Retry, then impute by cell-relative carry-forward}: $\hat p_{i,t} = p_{i,t-1}\times I_g^{t-1,t}$, i.e.\ move the missing item by its own cell's observed change. Flag as imputed. \\
\addlinespace[2pt]
Flight not operating (schedule change, cancellation) & Absent from schedule reference & \textbf{Remove from the matched sample} prospectively. A flight that does not exist has no price. Do not impute. \\
\addlinespace[2pt]
Sold out / fare class closed & Present in schedule, absent from availability & \textbf{This is a price signal, not missing data.} See \S\ref{sec:avail}. \\
\bottomrule
\end{tabularx}
\end{center}

\section{Stage 1: elementary index}

For each cell $g=(r,c,\tau)$ with matched quote set $\mathcal{M}_g$ present in both $0$
and $t$:
\begin{equation}
\boxed{\;
I_{g,t} \;=\; 100 \times \exp\!\left(\frac{1}{|\mathcal{M}_g|}\sum_{i\in\mathcal{M}_g}
\ln\frac{\tilde p_{i,t}}{\tilde p_{i,0}}\right)\;}
\end{equation}
Suppress and mark the cell if $|\mathcal{M}_g| < N_{\min}$ (default $N_{\min}=5$, derived
from the power calculation of \S3.8.3). A suppressed cell's weight is redistributed
pro-rata to the remaining cells within the same route --- exactly the mechanism MoSPI uses
when an item is unavailable in a market.

\section{Stage 2: carrier aggregation within a route}
\label{sec:carrierweights}

\begin{equation}
I_{r,\tau,t} \;=\; \sum_{c} \phi_{c|r}\; I_{(r,c,\tau),t}
\end{equation}
with the within-route carrier weight
\begin{equation}
\phi_{c|r} \;=\; \frac{S_{c}\cdot F_{c,r}\cdot \text{Seats}_{c,r}\cdot \text{PLF}_{c}}
{\sum_{c'} S_{c'}\cdot F_{c',r}\cdot \text{Seats}_{c',r}\cdot \text{PLF}_{c'}}
\end{equation}
where $F_{c,r}$ is weekly frequency on the route, $\text{Seats}_{c,r}$ the average seats per
departure, $\text{PLF}_c$ the carrier's published passenger load factor, and $S_c$ an
optional national-share prior for smoothing thin routes. In practice
$F\times\text{Seats}\times\text{PLF}$ is an estimate of \emph{passengers actually carried by
carrier $c$ on route $r$}, which is the economically correct weight.

\begin{keybox}
\bhead{Why this beats equal weighting.} IndiGo carries roughly two thirds of Indian domestic
passengers (67.4\% in July 2026). Under equal carrier weighting, a fare move at SpiceJet
(1.6\% share) would move the index as much as one at IndiGo. Under $\phi_{c|r}$, it moves it
proportionally to the passengers actually exposed to it. Load factors are published by DGCA
monthly (Akasa 92.2\%, SpiceJet 87.8\%, AI Group 85.5\%, IndiGo 85.1\% in June 2026), so
this weight is fully sourced.
\end{keybox}

\section{Stage 3: route aggregation}

\begin{equation}
I_{\tau,t} \;=\; \sum_{r\in\mathcal{B}} w_r\, I_{r,\tau,t},
\qquad
w_r \;=\; \frac{Q_r}{\sum_{r'\in\mathcal{B}} Q_{r'}}
\end{equation}
where $\mathcal{B}$ is the route basket and $Q_r$ is annual passengers on route $r$ from
DGCA/AAI sector-level traffic. Where the PSD supplies its own route list and weights ---
which is precisely what the PS means by \emph{``an index-construction module based on PSD
given routes and weights''} --- the module accepts them as a signed input file and uses
them in place of the derived weights. \textbf{Build the weight vector as an injectable
parameter from day one}; the whole point is that PSD, not the team, owns the weights.

\section{Stage 4: lead-time aggregation and the headline}

\begin{equation}
\boxed{\;\text{APIx}_t \;=\; \sum_{\tau\in\mathcal{T}} \omega_\tau\, I_{\tau,t}\;}
\end{equation}

\begin{warnbox}
\bhead{$\omega_\tau$ is the honest weak point of this project, and you should say so.}
The weights over advance-purchase windows should be the \emph{distribution of actual booking
lead times of Indian domestic travellers}. That distribution is held by the airlines and the
OTAs and is not published. Note that the DGCA asked airlines in December 2024 for exactly
this kind of ticket-level data --- booking date, fare charged, base fare --- and the
Federation of Indian Airlines refused on commercial-confidentiality grounds. So the data
exists and is contested; it is not something a hackathon team can obtain.

\textbf{APIx's answer:} treat $\omega_\tau$ as a declared, versioned policy parameter, ship
three named presets (\emph{Uniform}, \emph{Leisure-tilted}, \emph{Business-tilted}), publish
the index under all three, and include a tornado sensitivity chart showing how much the
headline moves across the range. Then state that MoSPI or DGCA can replace it with a
measured distribution the moment one exists. Turning your biggest gap into a published
sensitivity analysis is a strength move, not a weakness.
\end{warnbox}

\section{The dual time basis}
\label{sec:dual}

This is APIx's most distinctive methodological contribution, and it is directly grounded in
international practice.

\begin{center}
\small
\begin{tabularx}{\textwidth}{@{}L{3.6cm} X@{}}
\toprule
\bhead{Basis} & \bhead{Definition and rationale} \\
\midrule
\textbf{Booking basis} \newline (acquisition approach) & The quote is attributed to the
\emph{collection date} $t$. Answers: \emph{``what does it cost to buy a flight today?''}
This is the consumer-facing, regulator-facing, RBI-nowcast-facing series. \\
\addlinespace[3pt]
\textbf{Travel basis} \newline (use approach) & The quote is attributed to the
\emph{departure date} $d = t+\tau$. Answers: \emph{``what is the price of air travel
consumed in this month?''} This is the CPI-consistent series. The UK ONS states explicitly
that changes in air fares are recorded in the CPI \textbf{in the month in which the flight
departs, not when the ticket is bought.} \\
\bottomrule
\end{tabularx}
\end{center}

Formally, the travel-basis index at month $m$ aggregates over all quotes whose departure
falls in $m$, weighted by the lead-time distribution:
\begin{equation}
\text{APIx}^{\text{travel}}_m \;=\; \sum_{\tau\in\mathcal{T}} \omega_\tau \,
I_{\tau}\big(t = m - \tau\big)
\end{equation}
whereas the booking-basis index is
\begin{equation}
\text{APIx}^{\text{book}}_m \;=\; \sum_{\tau\in\mathcal{T}} \omega_\tau \, I_{\tau}(t=m)
\end{equation}

\noindent The two differ by a \emph{lead-time-weighted lag structure}. Publishing both and
showing the reconciliation is a genuinely novel contribution for an Indian price series, and
it directly demonstrates that the team read foreign NSI methodology rather than reinventing
one badly.

\section{Availability adjustment}
\label{sec:avail}

\subsection{The bias}

Suppose a cell normally offers fare families at Rs.\,4{,}000, Rs.\,6{,}000 and
Rs.\,9{,}000. During a demand surge the Rs.\,4{,}000 and Rs.\,6{,}000 buckets
sell out. A naive matched index observes only the Rs.\,9{,}000 quote --- but that quote
has not \emph{changed price}, so a strictly matched index reports \textbf{zero inflation}
during the exact episode a consumer experiences as a fare explosion. This is not a hypothetical;
it is the mechanism by which any matched-model fare index understates surge periods, and it
is the single most defensible methodological criticism of a naive implementation.

\subsection{The APIx treatment}

Define the availability ratio for cell $g$ at $t$:
\begin{equation}
A_g^{(t)} \;=\; \frac{\text{\# fare families observed offered at } t}{\text{\# fare families in the cell's reference bundle}}
\;\in\;(0,1]
\end{equation}

Then compute a \textbf{lowest-available-fare (LAF) index} alongside the matched index. The
LAF index tracks $\min_k \tilde p_{i,t,k}$ within the flight, i.e.\ the cheapest seat a
consumer could actually buy:
\begin{equation}
I^{\text{LAF}}_{g,t} = 100\times\exp\!\left(\frac{1}{|\mathcal{M}_g|}\sum_{i}
\ln\frac{\min_k \tilde p_{i,t,k}}{\min_k \tilde p_{i,0,k}}\right)
\end{equation}

The \textbf{Availability-Adjusted APIx} is a geometric blend controlled by $A$:
\begin{equation}
\boxed{\;
I^{\text{adj}}_{g,t} \;=\; \big(I^{\text{matched}}_{g,t}\big)^{A_g^{(t)}}\cdot
\big(I^{\text{LAF}}_{g,t}\big)^{1-A_g^{(t)}}\;}
\end{equation}
When everything is available ($A=1$) the adjusted index equals the matched index and nothing
changes. As buckets close ($A\to 0$) it converges to the lowest-available-fare index, which
does capture the surge. The blend is continuous, monotone, and reduces exactly to standard
practice in the normal case --- which is what makes it defensible rather than ad hoc.

\begin{keybox}
\bhead{This is the single best demo in the whole project.} Pick a real surge --- a festival
weekend, or a disruption episode --- and show three lines on one chart: the naive matched
index (flat, wrong), the LAF index (spiking), and the availability-adjusted index
(spiking, correctly). Then show the $A_g^{(t)}$ series underneath as an area chart
collapsing toward zero. Ninety seconds, and it proves the team understands price
measurement rather than data plumbing.
\end{keybox}

\section{The lead-time elasticity surface}

\subsection{The model}

Estimate the fare surface by pooled log-linear regression with fixed effects:
\begin{equation}
\ln \tilde p_{i} \;=\; \alpha
\;+\; \beta_1 \ln \tau_i
\;+\; \beta_2 (\ln \tau_i)^2
\;+\; \gamma_{r} \;+\; \lambda_{c} \;+\; \theta_{\text{dow}(d_i)}
\;+\; \psi_{\text{month}(d_i)}
\;+\; \rho\,\text{Holiday}_i
\;+\; \varepsilon_i
\end{equation}
where $\gamma_r$, $\lambda_c$, $\theta$, $\psi$ are route, carrier, day-of-week and
month-of-departure fixed effects and $\text{Holiday}_i$ is a festival/long-weekend
indicator. Standard errors clustered by flight.

\subsection{The elasticity}

The \emph{lead-time elasticity of fare} is the derivative of log price with respect to log
lead time:
\begin{equation}
\boxed{\;\eta(\tau) \;=\; \frac{\partial \ln p}{\partial \ln \tau} \;=\; \beta_1 + 2\beta_2\ln\tau\;}
\end{equation}
Interpretation: $\eta(\tau) = -0.30$ at $\tau=7$ means a 10\% increase in booking lead time
around one week out is associated with a 3\% lower fare. Because $\beta_2$ enters, the
elasticity is allowed to \emph{vary with lead time} --- which matters, since fares typically
fall from $\tau=45$ to a trough somewhere around $\tau=21$--$30$ and then rise sharply
approaching departure. The trough is at
\begin{equation}
\tau^\star = \exp\!\left(-\frac{\beta_1}{2\beta_2}\right)
\end{equation}
which yields a directly publishable consumer statistic: \textbf{``the cheapest time to book
\rt{del}--\rt{bom} is $\tau^\star$ days out.''} That single number is the most
press-friendly output of the entire system.

\subsection{Fare-bucket detection}

Because airlines price from discrete RBDs, the empirical fare distribution within a
route--carrier is multi-modal. Fit a one-dimensional Gaussian mixture in log space:
\begin{equation}
f(y) \;=\; \sum_{j=1}^{J}\pi_j\,\mathcal{N}\!\left(y;\mu_j,\sigma_j^2\right),
\qquad \sum_j \pi_j = 1
\end{equation}
with $J$ chosen by BIC. The fitted $\{\mu_j\}$ are estimates of the \textbf{fare ladder} ---
the airline's actual bucket structure, recovered from public prices without any privileged
data. Tracking $\mu_j$ over time separates \emph{``the airline raised its fare ladder''}
(structural) from \emph{``more people bought so we are selling from a higher bucket''}
(compositional). Very few teams will attempt this and it is directly relevant to DGCA's
tariff-compliance mandate under Rule 135 of the Aircraft Rules, 1937, which requires
airlines to charge within their declared tariff bands.

\section{Nowcasting CPI Transport}

APIx is available daily; the CPI air fare item index is monthly, published around the 12th
of the following month. That gap is a forecasting opportunity.

Specify a mixed-frequency bridge. Let $\pi^{\text{CPI}}_m$ be the monthly y-o-y inflation of
the air fare item and $x_m$ the monthly average of the APIx travel-basis daily series:
\begin{equation}
\pi^{\text{CPI}}_m \;=\; \alpha \;+\; \sum_{j=0}^{J}\beta_j\, x_{m-j}
\;+\; \sum_{k=1}^{K}\rho_k\,\pi^{\text{CPI}}_{m-k}
\;+\; \delta\,\text{ATF}_{m}
\;+\; u_m
\end{equation}
with an ATF-price control (ATF is 35--40\% of airline operating cost, and PPAC publishes
prices). Report:
\begin{itemize}
\item Out-of-sample RMSE against a naive random-walk benchmark. If APIx does not beat a
random walk, say so --- a negative result honestly reported is more credible than an
unvalidated claim.
\item A Granger-causality test of APIx $\to$ CPI air fare.
\item Lead--lag cross-correlation at daily and monthly frequency.
\end{itemize}

\begin{warnbox}
Do not claim a nowcast of \emph{headline} CPI. Air fare is a small slice of a division that
is 8.796\% of combined CPI. Claim a nowcast of the \textbf{air fare item index} and, at
most, a contribution to the Transport division. Overclaiming here is the fastest way to lose
a statistician judge.
\end{warnbox}

\section{The full computation, end to end}

\begin{center}
\small
\begin{tabularx}{\textwidth}{@{}C{0.9cm} L{4.0cm} X@{}}
\toprule
\bhead{Step} & \bhead{Operation} & \bhead{Formula / rule} \\
\midrule
0a & Validity gates & Deterministic rejects \\
0b & Decompose & $P^{\text{base}} = P^{\text{total}} - \text{levies} - \text{GST} - \text{CF}$ \\
0c & Outlier flag & Tukey $3\times$IQR and Hampel $|z^{\text{rob}}|>3.5$ on logs \\
0d & Impute / winsorise & Cell-relative carry-forward; cap at fences \\
1 & Elementary index & Jevons over matched $\kappa$ within $g=(r,c,\tau)$ \\
1b & Availability adjust & $I^{\text{adj}} = (I^{\text{matched}})^{A}(I^{\text{LAF}})^{1-A}$ \\
1c & De-seasonalise DOW & 7-day centred geometric moving average \\
2 & Carrier aggregation & $\phi_{c|r}$ from frequency $\times$ seats $\times$ PLF \\
3 & Route aggregation & $w_r$ from DGCA sector traffic, or PSD-supplied \\
4 & Lead-time aggregation & $\omega_\tau$, three declared presets \\
5 & Rebase and link & $\times\,100$; linking factor to CPI 2024 base \\
6 & Publish & Daily, weekly (7-day GM), monthly (calendar-month GM) \\
7 & Uncertainty & Delta-method SE $+$ flight-block bootstrap CI \\
8 & Robustness & GEKS-Jevons and TPD variants; publish the spread \\
\bottomrule
\end{tabularx}
\end{center}

\noindent Monthly aggregation uses the \emph{geometric} mean of daily indices, consistent
with MoSPI's own use of geometric means for annual averages in the linking-factor
computation:
\begin{equation}
\text{APIx}_m \;=\; \left(\prod_{t\in m}\text{APIx}_t\right)^{1/|m|}
\end{equation}
